\documentclass[11pt]{article}
\usepackage[margin=1in]{geometry}
\usepackage[T1]{fontenc}
\usepackage[utf8]{inputenc}
\usepackage{lmodern}
\usepackage{setspace}
\usepackage{hyperref}
\usepackage{booktabs}
\usepackage{enumitem}
\usepackage{amsmath,amssymb}
\usepackage[numbers]{natbib}
\onehalfspacing

\title{De la personnalité juridique des robots à la personnalité numérique de l'intelligence artificielle : \\ statut, imputation et responsabilité publique, 2019-2026}
\author{Researcher Agent French}
\date{3 avril 2026}

\begin{document}
\maketitle

\begin{abstract}
Cet article reprend un mémoire de bachelor suisse de 2019 consacré à la taxation des robots et montre comment son intuition juridique principale doit être reformulée en 2026. Le mémoire concluait qu'une taxe sur les robots, conçus comme capital productif, était en règle générale inefficiente, tout en suggérant qu'une imposition du revenu généré par des systèmes automatisés supposerait à terme une forme de personnalité électronique. Or, entre 2019 et 2026, le paysage normatif a changé. Le droit europeen et suisse n'a pas retenu la voie d'une personnalité juridique générale de l'intelligence artificielle. Il a privilégié une régulation par les acteurs, les usages et les risques. L'article soutient que le mémoire avait correctement identifié le vrai problème doctrinal, mais qu'il l'orientait vers une solution trop lourde. La contribution la plus défendable en 2026 n'est pas la personnalité pleine et entière de l'IA, mais une personnalité numérique fonctionnelle et limitée, destinée à améliorer l'imputation, la traçabilité et certains effets fiscaux ou assurantiels dans des contextes strictement circonscrits. Une telle construction permettrait de raffiner l'architecture de responsabilité sans dessaisir les personnes physiques et morales de leur rôle central.
\end{abstract}

\noindent\textbf{Mots-clés :} intelligence artificielle, robots, personnalité juridique, personnalité numérique, fiscalité, responsabilité, Suisse

\section{Introduction}
Le débat sur la taxation des robots, particulièrement vif à la fin des années 2010, reposait sur une idée simple: si l'automatisation remplace une partie du travail humain, la base fiscale assise sur le revenu du travail et sur les cotisations sociales risque de s'éroder. Le mémoire de 2019 sur lequel se fonde le présent projet s'inscrivait pleinement dans ce cadre. Son objectif principal consistait à déterminer si la taxation des robots permettait, dans une petite économie ouverte comme la Suisse, de préserver les recettes publiques et de soutenir le financement des assurances sociales.

Ce mémoire aboutissait à une conclusion nuancée mais globalement sceptique. La taxe sur les robots y apparaissait comme une solution coûteuse, difficile à administrer, potentiellement défavorable à l'innovation, et finalement peu souhaitable dans la plupart des hypothèses ordinaires. Toutefois, au-delà de cette critique économique, le texte formulait une intuition doctrinale particulièrement féconde: dès que l'on cesse de considérer les robots comme de simples biens productifs pour leur attribuer un revenu ou une capacité contributive, la question de leur statut juridique devient inévitable.

C'est ici que le mémoire parlait de ``personnalité juridique électronique''. En 2026, cette intuition mérite d'être conservée, mais aussi profondément reformulée. Le débat n'est plus centré sur le robot matériel, individualisable, localisable dans une usine ou dans une chaîne de production. Il s'est déplacé vers des systèmes d'IA distribués, logiciels, intégrés à des services, ou déployés à travers des infrastructures transfrontières. Dans le même temps, les principales évolutions juridiques europeennes et suisses n'ont pas consacré la personnalité de l'IA. Elles ont préféré renforcer les obligations pesant sur les fournisseurs, les déployeurs et les utilisateurs de systèmes à risque.

L'hypothèse centrale de cet article est donc la suivante: le mémoire de 2019 avait bien repéré le bon problème, mais non la bonne solution finale. En 2026, l'option la plus défendable n'est pas une personnalité juridique générale de l'IA. C'est l'élaboration éventuelle d'une personnalité numérique fonctionnelle et limitée, conçue comme un outil d'imputation et de gouvernance dans des cas exceptionnels. Une telle figure ne créerait pas un nouveau sujet de droit comparable à une personne physique ou morale. Elle servirait plutôt à organiser plus rigoureusement la traçabilité, la documentation, et certains effets fiscaux ou assurantiels liés à l'usage de l'IA.

\section{Relecture du cadre de 2019}
Le mémoire de 2019 s'appuyait sur une littérature solide relative à la fiscalité des robots, à l'automatisation, à la part du travail dans le revenu national et au financement des assurances sociales. Sa démonstration économique reposait sur l'efficience fiscale, l'incidence fiscale et une modélisation de la production permettant d'évaluer les effets d'une taxe sur les robots à court et à long terme.

Plusieurs éléments de ce cadre conservent aujourd'hui leur pertinence. D'abord, la critique de l'objet imposable demeure convaincante. Une taxe assise sur le robot comme bien de production souffre d'une définition incertaine, d'une forte vulnérabilité à l'optimisation et d'une difficulté de localisation, surtout lorsque l'automatisation prend une forme logicielle. Ensuite, l'analyse de l'incidence garde une vraie force explicative: dans un contexte concurrentiel et ouvert, les charges nouvelles risquent d'être répercutées sur l'emploi, les salaires, la compétitivité ou l'investissement. Enfin, le mémoire avait raison d'identifier la personnalité juridique comme point de bascule dès lors qu'on cherche à dépasser une taxation des inputs pour raisonner sur une capacité contributive attribuée.

Cependant, le cadre de 2019 est désormais trop étroit pour saisir la réalité de l'IA en 2026. Il repose encore sur une représentation du progrès technique comme remplacement du travail par un capital largement identifiable. Or l'IA actuelle agit souvent sous forme de modèles generalistes, d'outils d'aide à la décision, de services intégrés et de systèmes en réseau. Le problème n'est donc plus simplement celui du robot comme machine, mais celui de l'IA comme systeme socio-technique inséré dans des chaînes d'acteurs.

\section{Le déplacement du problème juridique}
Le déplacement du débat entre 2019 et 2026 peut être formulé de la manière suivante: on est passé d'une interrogation sur le robot comme objet à une interrogation sur l'IA comme architecture de responsabilités.

En 2019, la personnalité juridique apparaissait surtout comme une hypothèse spéculative. Elle servait à résoudre une difficulté fiscale: si l'on veut taxer le revenu d'un robot, il faut lui reconnaître une forme de statut. En 2026, la question est plus large et plus pratique. Les systèmes d'IA peuvent affecter des droits fondamentaux, générer des contenus, orienter des décisions, restructurer le travail et influencer les relations entre entreprises, travailleurs et autorités publiques. Pourtant, les ordres juridiques n'ont pas choisi de répondre à ces défis en transformant l'IA en personne de droit.

Ce choix s'explique. La personnalité juridique n'est pas une récompense symbolique attribuée à des entités sophistiquées. C'est une technique normative d'imputation. Or, pour les besoins actuels du droit, les destinataires les plus adaptés des obligations demeurent les personnes physiques, les entreprises, les autorités publiques et les autres organisations capables d'être surveillées, sanctionnées, tenues d'informer, de prévenir les risques, de réparer les dommages et d'adapter leur conduite.

Ainsi, le débat contemporain invite à distinguer au moins trois niveaux. Le premier est celui de la personnalité juridique pleine et entière, qui paraît aujourd'hui excessif. Le second est celui d'une absence totale de statut propre, combinée à une régulation intégralement centrée sur les acteurs humains et organisationnels. Le troisième, plus subtil, est celui d'une personnalité numérique limitée, qui ne confère pas un ensemble complet de droits et d'obligations, mais permet de mieux ordonner l'imputation de certains effets juridiques.

\section{Pourquoi la personnalité pleine de l'IA reste peu défendable}
L'idée d'une personnalité juridique générale de l'IA demeure fragile pour plusieurs raisons.

D'abord, elle exigerait une refonte profonde de catégories fondamentales du droit privé et du droit public. Une personne juridique n'est pas seulement un point abstrait d'imputation. Elle suppose une capacité d'agir juridiquement, une représentation, un patrimoine, des mécanismes de responsabilité, voire une articulation avec des droits subjectifs. Accorder un tel statut à l'IA ouvrirait immédiatement des questions de cohérence sur la propriété, la responsabilité civile, la faillite, la sanction, la fiscalité et même les libertés fondamentales.

Ensuite, une telle construction risquerait de produire un effet de diversion. Au lieu d'accroître la responsabilité des opérateurs humains et des entreprises, elle pourrait devenir un écran conceptuel derrière lequel les véritables centres de pouvoir se dérobent. Le droit perdrait alors en lisibilité ce qu'il gagnerait peut-être en sophistication.

Enfin, la personnalité pleine de l'IA paraît socialement et politiquement mal justifiée. Les principaux enjeux contemporains - transparence, supervision humaine, sécurité, non-discrimination, documentation, recours - peuvent être traités sans faire de l'IA un sujet de droit autonome.

\section{Vers une personnalité numérique fonctionnelle et limitée}
Le rejet de la personnalité générale ne conduit pas nécessairement au rejet de toute innovation statutaire. Il ouvre au contraire un espace plus précis: celui d'une personnalité numérique fonctionnelle et limitée.

Une telle figure ne vise pas à humaniser la machine, ni à reproduire par analogie la personnalité morale des sociétés. Elle poursuit un objectif strictement instrumental. Elle sert à identifier certains systèmes ou certains déploiements d'IA comme points stabilisés d'imputation documentaire et organisationnelle.

Pour être juridiquement défendable, cette personnalité numérique devrait respecter cinq conditions.

Premièrement, elle devrait etre \emph{strictement fonctionnelle}. Sa raison d'être ne serait ni philosophique ni symbolique, mais pratique.

Deuxièmement, elle devrait etre \emph{limitée à des contextes définis}. Tous les systèmes d'IA n'ont pas besoin d'un tel statut. La catégorie ne devrait concerner que certains usages où les mécanismes ordinaires d'imputation apparaissent insuffisamment précis.

Troisièmement, elle devrait etre \emph{subordonnée aux responsabilités humaines et organisationnelles}. L'existence d'une personnalité numérique ne devrait jamais permettre à une entreprise ou à un opérateur de se décharger de ses obligations.

Quatrièmement, elle devrait etre \emph{favorable à la traçabilité}. Son utilité normative principale serait d'améliorer l'identification des systèmes concernés, leur historique de déploiement, leurs journaux, les personnes responsables et les obligations attachées à leur usage.

Cinquièmement, elle devrait etre \emph{compatible avec le droit existant}. Il ne s'agirait pas de créer un ordre juridique parallèle pour l'IA, mais de compléter ponctuellement les catégories disponibles.

Dans cette perspective, la personnalité numérique pourrait servir d'interface juridique dans certains environnements fortement réglementés. Elle pourrait faciliter l'enregistrement, la documentation, la désignation d'un responsable, la traçabilité des versions et, dans des cas circonscrits, l'articulation de certains effets fiscaux ou assurantiels. Mais le centre de gravité normatif resterait du côté des personnes physiques et morales.

\section{Conséquences fiscales et assurantielles}
La force du mémoire initial était de montrer que la fiscalité des robots ne pouvait être pensée indépendamment du financement des assurances sociales. Cette intuition conserve toute sa valeur. La question fiscale de 2026 n'est plus de savoir s'il faut taxer un robot matériel. Elle consiste à déterminer comment le droit fiscal et social peut continuer à suivre une création de valeur de plus en plus liée à des actifs immatériels, à des services d'IA et à des chaînes d'acteurs complexes.

Dans ce contexte, la personnalité numérique limitée ne devrait pas être comprise comme la création d'un nouveau contribuable autonome de portée générale. Une telle solution réintroduirait les difficultés déjà signalées en 2019: double imposition, définition du revenu attribuable, problème de la localisation, faible transparence de l'objet imposable.

Sa fonction serait plus modeste. Elle pourrait contribuer à clarifier l'imputation de certains effets et à stabiliser la chaîne documentaire nécessaire au traitement fiscal ou assurantiel de certains usages d'IA. En d'autres termes, la personnalité numérique n'aurait pas vocation à remplacer le contribuable humain ou organisationnel. Elle aiderait à mieux relier l'événement technique aux responsabilités juridiques et économiques qui en découlent.

L'apport le plus solide du mémoire de 2019 apparaît alors sous un jour nouveau. Il ne s'agit pas d'avoir démontré qu'il fallait créer une personnalité électronique générale, mais d'avoir compris que toute tentative de faire porter à l'automatisation des conséquences fiscales directes soulève une question d'imputation juridique. En 2026, cette question demeure entière, mais la réponse la plus robuste passe par une construction limitée.

\section{Objections et précautions}
Une première objection consiste à dire que la personnalité numérique est superflue. Peut-être le droit positif, enrichi par des obligations pesant sur les fournisseurs, les déployeurs et les autorités, suffit-il déjà. Cette objection doit être prise au sérieux. La personnalité numérique ne se justifie que si elle apporte un gain réel de clarté normative.

Une deuxième objection porte sur le risque d'anthropomorphisme. Le vocabulaire de la personnalité peut entretenir des confusions entre capacité technique, autonomie sociale et statut juridique. Pour cette raison, toute réforme éventuelle devrait expliciter qu'aucun droit subjectif général n'est attaché à la personnalité numérique.

Une troisième objection vise le risque de déresponsabilisation des entreprises. Si la catégorie servait de support à une externalisation de la faute ou du risque, elle deviendrait contre-productive. Il faut donc que le droit prévoie expressément que la personnalité numérique ne fait jamais écran à la responsabilité des acteurs humains et organisationnels.

Enfin, une dernière objection est comparative. La Suisse, comme d'autres États européens, privilégie aujourd'hui une combinaison de mesures contraignantes ciblées et de mesures non contraignantes. Dans ce contexte, la personnalité numérique doit rester une hypothèse doctrinale et non un programme legislatif présumé inévitable.

\section{Conclusion}
Le mémoire de 2019 sur la taxation des robots conserve une valeur remarquable, non pas parce qu'il aurait apporté la solution définitive à la question fiscale de l'automatisation, mais parce qu'il a vu plus loin que le débat immédiat sur la taxe. Il a compris que la difficulté réelle se situait du côté du statut juridique de l'entité technologique dès lors qu'on voulait lui attribuer un revenu, une capacité contributive ou une fonction quasi-travaillante.

En 2026, cette intuition doit être conservée mais disciplinée. Le droit positif suisse et europeen ne s'oriente pas vers une personnalité juridique générale de l'IA. Il organise plutôt des obligations portant sur les fournisseurs, les déployeurs, les utilisateurs et les contextes à risque. Cette orientation est, dans l'état actuel des techniques et des institutions, largement justifiée.

Pour autant, la question statutaire n'a pas disparu. Elle s'est resserrée. Ce qui peut aujourd'hui être défendu n'est pas la personnification complète de l'IA, mais l'examen d'une personnalité numérique fonctionnelle et limitée, destinée à améliorer l'imputation, la traçabilité et certains effets juridiques spécialisés. Une telle solution permettrait de reprendre la meilleure intuition du mémoire de 2019 tout en évitant l'erreur la plus probable du débat contemporain: croire que l'IA doit devenir une personne pour être convenablement gouvernée.

\bibliographystyle{plainnat}
\begin{thebibliography}{99}
\bibitem{memoire2019} Graz, David. \textit{Taxation des robots}. Mémoire de bachelor, 2019.
\bibitem{oberson2017} Oberson, Xavier. \textit{Taxer les robots ? L'émergence d'une capacité contributive électronique}. 2017.
\bibitem{floridi2017} Floridi, Luciano. \textit{Robots, Jobs, Taxes, and Responsibilities}. 2017.
\bibitem{acemoglu2018} Acemoglu, Daron et Restrepo, Pascual. \textit{Artificial Intelligence, Automation and Work}. 2018.
\bibitem{guerreiro2019} Guerreiro, Joao, Rebelo, Sergio et Teles, Pedro. \textit{Should Robots Be Taxed?} 2019.
\bibitem{thuemmel2018} Thuemmel, Uwe. \textit{Optimal Taxation of Robots}. 2018.
\bibitem{euia2026} Documents de la Commission europeenne sur l'AI Act, consultés en 2026.
\bibitem{suisse2025} Documents du Conseil federal suisse sur la réglementation de l'IA, 2025-2026.
\bibitem{coe2024} Convention-cadre du Conseil de l'Europe sur l'intelligence artificielle, 2024.
\end{thebibliography}

\end{document}
